%% kapa: exam=1 round=31 year=2020 period=2 form=A subject=appraisal_law subject_ko=감정평가관계법규 seq=6 number=6
%% source: Q-Net 감정평가사 기출문제 (cst00309) artlSeq=5207219 file=extracted/1008/5207219/A형/A형/2020년도 제31회 감정평가사 1차 시험_2교시_A형.pdf
%% conversion: kordoc 4.24.51 -> tools/gen1.py (auto); review=proofread-sonnet; flags=none
\begin{question}{6}
국토의 계획 및 이용에 관한 법령상 도시ㆍ군계획 등에 관한 설명으로 옳은 것은?
\begin{choices}
  \item 광역도시계획은 광역계획권의 장기발전방향을 제시하는 계획을 말한다.
  \item 도시ㆍ군기본계획의 내용이 광역도시계획의 내용과 다를 때에는 도시ㆍ군기본계획의 내용이 우선한다.
  \item 도시ㆍ군관리계획으로 결정하여야 할 사항은 국가계획에 포함될 수 없다.
  \item 시장 또는 군수가 관할 구역에 대하여 다른 법률에 따른 환경에 관한 부문별 계획을 수립할 때에는 도시ㆍ군관리계획의 내용에 부합되게 하여야 한다.
  \item 이해관계자가 도시ㆍ군관리계획의 입안을 제안한 경우, 그 입안 및 결정에 필요한 비용의 전부를 이해관계자가 부담하여야 한다.
\end{choices}
\end{question}
\answer{1}
